\section{基于运行生产资料的汽轮机系统故障知识图谱构建}
\label{sec:knowledge_graph_v3}

\subsection{引言}

燃煤机组汽轮机系统由多级汽缸、再热系统、抽汽回热系统、给水系统及凝汽器冷端等多个设备共同组成。不同设备之间通过蒸汽、凝结水和给水形成连续的物质与能量传递关系，其运行状态并非相互独立。当局部设备发生通流、汽封、阀门或换热性能异常时，故障影响往往首先改变故障设备内部压力、温度和流量分布，随后通过排汽、抽汽、再热或冷端接口向相邻设备传播，最终在多个监测参数上形成具有空间关联的异常响应。第二章建立的状态监测模型能够给出各状态参数在当前运行边界下的健康参考，但参数偏差本身并不能直接表示“异常参数属于哪个设备”“不同异常之间是否存在实际热力联系”以及“多个异常共同指向何种故障机理”。因此，需要进一步利用汽轮机系统长期积累的设备结构、运行经验和故障知识，对不同参数之间的工程关联进行结构化表达。

知识图谱（Knowledge Graph, KG）以实体和关系为基本单元，将分散于不同资料中的知识组织为图结构\cite{Hogan2021KnowledgeGraphs}。相较于仅依据相关系数或数据相似度建立变量连接，领域知识图谱能够显式表示设备归属、介质流向、故障因果和异常表现等工程关系，使后续诊断模型不仅利用参数数值变化，还能够利用参数之间的物理连接关系。已有汽轮机图神经网络研究表明，汽封缺陷、蒸汽泄漏和汽耗增加等故障信息可依据设备结构形成分层传播图，并利用图网络完成故障传播关系推断\cite{ZhangHu2021FaultPropagationGNN}。因此，将汽轮机运行生产资料中的设备与故障知识组织为领域知识图谱，可以为下一章KG-GCN模型提供具有工程含义的静态拓扑先验。

本章借鉴杨昊东在CFB锅炉主辅机系统中利用运行生产资料构建故障知识图谱的总体思路\cite{YangHaodong2026CFB}，结合汽轮机系统的主通流、再热、抽汽回热及冷端特征重新定义实体类型和关系类型。首先分析汽轮机系统典型热力故障及其参数响应，明确图谱需要描述的知识范围；随后从系统结构资料、DCS测点信息、设备运行与检修知识以及公开故障文献中提取设备、参数、异常、故障和处置等实体，构建候选知识三元组；在此基础上通过实体类型规则、关系方向规则、语义相似度和证据来源对候选知识进行校核与归一化，形成汽轮机故障知识图谱核心语义层。最后建立领域参数节点与实际DCS测点之间的映射，为下一章将动态状态偏差加载到图节点提供基础。

\subsection{燃煤机组汽轮机系统参数及故障特性}
\label{subsec:turbine_fault_characteristics_v3}

\subsubsection{汽轮机系统参数关联与故障传播}

汽轮机系统主通流沿“主蒸汽--VHP--一次再热RH1--HP--二次再热RH2--IP--LP--凝汽器”的方向传递。主蒸汽压力、温度和流量首先决定VHP入口热力状态，VHP排汽参数又构成一次再热系统入口条件；一次热再参数进一步成为HP入口边界，HP排汽经二次再热后形成IP入口状态，IP和LP最终受凝汽器背压约束完成低压侧膨胀过程。除主通流外，汽轮机各级抽汽向高、低压加热器及除氧器提供热源，凝结水经低压加热、除氧、给水泵升压和高压加热后返回锅炉，凝汽器则通过循环水系统形成冷端换热边界。由此形成的主通流、抽汽回热、给水和冷端回路共同构成汽轮机系统的参数关联网络。

在正常运行情况下，设备入口边界、内部状态和出口状态之间存在相对稳定的热力映射。例如，VHP主蒸汽边界与叶片级压力、排汽温度及一级抽汽状态具有确定的工况相关性；LP末级状态与凝汽器压力之间存在直接边界联系；抽汽参数又同时影响对应加热器的换热过程。当设备发生故障时，原有映射被破坏，异常会沿上述结构关系形成多参数协同变化。以汽封泄漏为例，汽封间隙增大首先改变局部有效蒸汽流量，进一步影响汽缸效率和排汽状态；若泄漏发生在低压轴端，还可能伴随空气漏入并影响凝汽器真空。Zhang和Hu在汽轮机故障传播研究中将汽封缺陷、蒸汽泄漏和汽耗变化组织成层次化故障传播图，说明故障关系具有明显的结构传播特征\cite{ZhangHu2021FaultPropagationGNN}。

因此，本文知识图谱不仅表示“故障与某个参数相关”，还同时保存三类知识：第一类为设备和参数之间的静态结构关系，用于确定参数归属和介质传播方向；第二类为故障、异常和参数之间的因果与表现关系，用于描述故障发生后状态如何变化；第三类为故障与检修处置之间的经验关系，用于形成从异常识别到故障解释和处理建议的完整知识链。上述结构为后续在具有物理联系的参数之间开展图卷积提供依据。

\subsubsection{典型热力故障及参数响应特征}

本文面向DCS热工参数能够反映的通流与热力故障开展知识构建，不将转子不平衡、轴承振动等需要专用振动信号分析的机械故障作为主要对象。结合汽轮机故障文献及系统可观测参数，重点整理调节阀异常、汽封泄漏、通流部件侵蚀与沉积、凝汽器换热性能劣化以及回热系统泄漏等典型故障。

（1）\textbf{调节阀死区及配汽异常。} 高压调节阀通过改变有效流通面积调节进入汽轮机的蒸汽流量。执行机构磨损、卡涩或机械间隙可能形成阀门死区，即控制指令变化而实际阀位在一定范围内不产生相应动作。Zhang和Hu针对汽轮机控制阀建立图模型时指出，阀门死区会使指令与实际位置关系发生异常，并可能诱发高压侧及系统振荡\cite{ZhangHu2021ControlValveGraph}。因此，可将“高压调节阀死区--导致--阀位对指令不敏感--表现为--实际阀位异常”作为一条典型故障链，同时将主蒸汽流量和入口压力波动作为关联状态。

（2）\textbf{汽封磨损与蒸汽泄漏。} 级间汽封和轴端汽封用于限制不同压力区域之间的蒸汽泄漏，汽封齿磨损或间隙增大会增加无效泄漏量。汽轮机故障传播研究表明，HP、IP和LP的级间或轴端汽封缺陷可导致蒸汽泄漏并进一步引起汽耗增加；低压轴端密封恶化还可能导致空气漏入低压系统，造成凝汽器真空下降\cite{ZhangHu2021FaultPropagationGNN}。针对HP--IP中间汽封泄漏的热力研究进一步表明，中间汽封泄漏与IP效率下降密切相关，热再温度和IP排汽温度是泄漏诊断的重要测量参数\cite{Xu2011HPIPLeakage}。因此，汽封故障适合表示为“部件缺陷--泄漏异常--效率/汽耗变化--压力温度响应”的多层关系。

（3）\textbf{通流部件侵蚀与沉积。} 汽轮机长期运行后，叶片、喷嘴和汽封等通流部件可能出现固体颗粒侵蚀、表面粗糙度增加、沉积及间隙变化。Kubiak等基于汽轮机现场诊断与检修结果指出，固体颗粒侵蚀会改变喷嘴有效面积并引起关键压力分布变化，叶片表面沉积则会减小有效通流面积并造成效率和出力下降\cite{Kubiak2002TurbineDegradation,Kubiak2002ScaleDeposit}。这类故障通常无法由单一DCS量直接测量，但其引起的通流面积、压力分布和排汽状态变化能够在多个热工参数上共同表现，因而特别适合利用知识图谱描述多参数之间的诊断关联。

（4）\textbf{凝汽器污垢及冷端性能劣化。} 凝汽器水侧污垢会增加换热热阻，降低凝汽器传热能力。针对燃煤机组的在线污垢监测研究表明，污垢可通过换热有效性或污垢热阻变化进行表征\cite{Li2016CondenserFouling}；凝汽压力同时受到机组负荷、循环水入口温度、循环水流量和换热面状态影响，凝汽压力升高会降低低压缸可利用焓降并影响机组效率\cite{Medica2018CondenserCondition}。因此，图谱中需要同时表示设备故障和外部冷端边界，形成“凝汽器水侧污垢--换热能力下降--凝汽压力升高--LP有效焓降降低”的知识链，以便后续模型区分环境边界变化与设备性能异常。

（5）\textbf{给水加热器内部泄漏及回热性能下降。} 给水加热器内部管束或隔板泄漏会改变正常抽汽--给水换热过程，使换热量和出口给水温度发生变化。Barszcz和Czop建立的给水加热器模型表明，换热功率等过程参数可用于跟踪内部泄漏等慢性性能变化\cite{Barszcz2011FeedwaterHeater}；Heo和Lee进一步研究了闭式给水加热器内部泄漏的识别与定位\cite{Heo2012FWHLeakage}。据此，图谱中可形成“加热器内部泄漏--回热性能下降--机组热效率下降”的故障链，并将抽汽侧压力、温度及给水出口温度作为关联监测参数。

\begin{table}[htbp]
    \centering
    \caption{汽轮机系统典型故障及知识图谱表达}
    \label{tab:turbine_fault_knowledge_v3}
    \begin{tabular}{p{0.20\textwidth}p{0.27\textwidth}p{0.28\textwidth}p{0.17\textwidth}}
        \toprule
        故障类型 & 主要异常机理 & 关联参数/状态 & 主要系统影响 \\
        \midrule
        调节阀死区 & 指令变化不能及时形成有效阀位响应 & 阀位反馈、主汽流量、入口压力波动 & 配汽与负荷响应异常 \\
        汽封磨损/泄漏 & 有效蒸汽流量减少或空气漏入 & 排汽状态、效率、汽耗、凝汽器真空 & 汽缸效率及冷端状态变化 \\
        通流侵蚀 & 喷嘴/叶片有效面积与粗糙度变化 & 级组压力、排汽参数、效率 & 通流能力和出力变化 \\
        叶片沉积 & 有效通流面积减小 & 级前压力、效率、出力 & 性能退化 \\
        凝汽器污垢 & 换热热阻增加 & 背压、真空、循环水温度与流量 & LP焓降和机组效率下降 \\
        加热器内漏 & 抽汽--给水换热过程被破坏 & 抽汽侧参数、给水出口温度 & 回热性能下降 \\
        \bottomrule
    \end{tabular}
\end{table}

\subsection{基于运行生产资料的知识图谱构建方法}
\label{subsec:kg_construction_v3}

\subsubsection{汽轮机系统运行生产资料}

电厂长期运行过程中形成的故障知识并不集中存在于单一数据库中，而是分布在设计说明、系统图、DCS测点表、运行规程、检修记录、缺陷记录和故障分析报告等不同资料中。不同资料所提供的信息粒度和作用并不相同。设备设计资料和系统图重点描述设备结构、介质流向和设计边界，可用于确定“设备--设备”和“参数--设备”的稳定结构关系；DCS测点表给出现场可观测参数名称、点号和单位，是将领域知识连接到实时运行数据的基础；运行规程、检修记录和故障分析资料通常包含故障现象、参数变化、原因判断和处理过程，可用于提取“故障--异常--参数--处置”关系；同行评议文献则用于补充和交叉验证典型故障机理。

本文当前核心语义层主要使用已经核验的系统拓扑、DCS参数定义以及同行评议文献中的故障机理关系进行构建。对于后续能够获得的大规模值长日志、缺陷记录和检修文本，采用同一套知识抽取流程进行扩展。科研智能体在该流程中主要承担文献检索、候选关系定位、术语归并建议和候选三元组生成，候选结果仍需经过实体类型、关系方向和证据来源校验后方可进入核心图谱。

\begin{table}[htbp]
    \centering
    \caption{汽轮机故障知识来源及其主要用途}
    \label{tab:kg_sources_v3}
    \begin{tabular}{p{0.22\textwidth}p{0.30\textwidth}p{0.36\textwidth}}
        \toprule
        知识来源 & 典型内容 & 图谱构建用途 \\
        \midrule
        系统图及设备资料 & 主通流、抽汽、回热、冷端连接关系 & 设备拓扑、介质流向、设备归属 \\
        DCS测点表 & 点号、中文名称、单位、测量位置 & 参数实体、参数--设备归属、测量映射 \\
        运行/检修/缺陷资料 & 故障现象、参数异常、原因和处理措施 & 故障、异常、处置关系 \\
        同行评议文献 & 汽封泄漏、通流退化、凝汽器污垢等机理 & 故障机理补充及交叉验证 \\
        \bottomrule
    \end{tabular}
\end{table}

\subsubsection{节点与关系的构建及提取}

根据汽轮机故障诊断任务对知识的需求，本文定义设备、参数、异常、故障和处置五类实体。设备实体用于表示VHP、HP、IP、LP、再热器、凝汽器、加热器等物理对象；参数实体对应能够由DCS测量或由运行数据计算得到的压力、温度、流量、阀位和效率等状态；异常实体描述参数或性能发生偏离后的中间状态，例如“蒸汽泄漏”“凝汽器真空下降”和“通流面积减小”；故障实体表示具体故障模式；处置实体则保存检修或运行处理措施。

设领域实体集合为
\begin{equation}
    V=V_D\cup V_P\cup V_A\cup V_F\cup V_O,
\end{equation}
其中$V_D$、$V_P$、$V_A$、$V_F$和$V_O$分别表示设备、参数、异常、故障和处置实体集合。关系集合定义为
\begin{equation}
    \mathcal R=\{r_{bel},r_{flow},r_{occur},r_{cause},r_{affect},r_{show},r_{monitor},r_{action}\},
\end{equation}
分别表示归属、介质/能量连接、发生于、导致、影响、表现为、关联监测和处置关系。由此，一个知识事实可表示为三元组
\begin{equation}
    \tau=(h,r,t),\qquad h,t\in V,\ r\in\mathcal R.
\end{equation}
例如，“凝汽器水侧污垢--导致--凝汽器换热能力下降”“凝汽器换热能力下降--影响--凝汽器压力/背压”“VHP排汽温度--归属--VHP”等均可表示为统一三元组形式。

对于文本类资料，首先利用大语言模型按照预先规定的实体类型和关系类型生成候选三元组。为降低自由生成造成的关系漂移，提示模板限定允许输出的实体类别、关系类别和固定JSON/表格格式，并要求同时返回支持该关系的原始句段。随后采用领域规则进行第一次校验：设备节点不能作为数值参数处理；“发生于”关系的尾实体应为设备；“处置为”关系应由故障或异常指向处置；参数归属关系应由参数指向设备；介质流向须符合汽轮机实际热力方向。未满足schema约束的候选关系直接剔除或进入人工复核列表。

不同资料中同一实体可能具有多种表述，例如“凝汽器背压”“凝汽压力”和“低压缸排汽背压”在不同场景下可能对应相同或高度相关的领域概念。为减少同义实体造成的节点重复，采用字符串标准化与句向量语义相似度进行实体对齐。设候选实体$a$与标准实体$b$的语义相似度为
\begin{equation}
    s_{sem}(a,b)=\cos(\mathbf z_a,\mathbf z_b),
\end{equation}
其中$\mathbf z_a$和$\mathbf z_b$为Sentence-BERT等句向量模型得到的语义表示\cite{Reimers2019SentenceBERT}。对于高相似度候选，再结合设备归属、单位和上下文进行二次校验，以避免仅凭名称相似将不同测点错误合并。

在证据校核阶段，每条候选关系均保存来源标识。本文将实际系统结构、DCS点表和直接支持关系的原始同行评议文献作为高优先级证据，将跨机型通用故障机理或间接支持关系作为辅助证据。只有能够通过系统结构兼容性和来源证据审核的关系才进入核心语义图；其余关系保存在候选库中供后续运行资料补充验证。该处理既保留知识图谱的可扩展性，也避免将单次模型生成结果直接作为确定的工程知识。

\begin{figure}[htbp]
    \centering
    \fbox{\parbox[c][5.2cm][c]{0.90\textwidth}{\centering
    图3-1占位：系统图/DCS点表/运行与检修资料/故障文献\\
    $\Downarrow$ 科研智能体与LLM生成候选三元组\\
    $\Downarrow$ 实体类型与关系方向规则校验\\
    $\Downarrow$ Embedding实体标准化与来源证据审核\\
    $\Downarrow$ 汽轮机故障知识图谱核心语义层。}}
    \caption{汽轮机系统故障知识图谱构建流程}
    \label{fig:kg_construction_flow_v3}
\end{figure}

\subsection{汽轮机系统故障知识图谱构建结果}
\label{subsec:kg_results_v3}

依据上述实体和关系模式，结合汽轮机系统拓扑、DCS参数概念及故障文献证据，构建了第一版汽轮机系统故障知识图谱核心语义层。该语义层共包含107个节点和114条三元组，其中设备节点23个、参数节点39个、异常节点24个、故障节点15个、处置节点6个。关系边主要包括参数归属36条、故障/异常导致关系24条、设备介质或能量连接15条、故障发生于设备10条、关联监测8条以及处置关系6条，此外还包含主通流介质流向、影响和表现关系。

\begin{table}[htbp]
    \centering
    \caption{汽轮机故障知识图谱核心语义层统计}
    \label{tab:kg_stats_v3}
    \begin{tabular}{lrrrrr}
        \toprule
        类型 & 设备 & 参数 & 异常 & 故障 & 处置 \\
        \midrule
        节点数量 & 23 & 39 & 24 & 15 & 6 \\
        \bottomrule
    \end{tabular}
\end{table}

核心语义层能够同时表达系统结构知识和故障机理知识。例如，主通流中保存“VHP--介质流向--一次再热系统RH1”等设备连接关系；参数层保存“VHP排汽温度--归属--VHP”等测量归属关系；故障层保存“固体颗粒侵蚀--导致--首级喷嘴面积增大”“凝汽器水侧污垢--导致--凝汽器换热能力下降”等故障机理关系。通过故障、异常和参数节点之间的连接，可从一个故障实体追溯至其可能影响的参数，也可以从异常参数反向检索与其具有知识关联的设备和故障实体。

\begin{figure}[htbp]
    \centering
    \fbox{\parbox[c][5.2cm][c]{0.90\textwidth}{\centering
    图3-2占位：展示VHP、RH1、HP、IP、LP、凝汽器等设备节点；\\
    参数节点环绕设备分布；故障与异常节点连接对应参数；\\
    用不同线型区分“归属/介质连接/导致/关联监测”等关系。}}
    \caption{汽轮机系统故障知识图谱核心语义层示意图}
    \label{fig:kg_core_graph_v3}
\end{figure}

知识图谱中的参数节点属于领域概念层，实际诊断所使用的数据来自具体DCS测点。因此，在核心语义层之外建立参数概念与实际测点之间的测量映射
\begin{equation}
    \mathcal M:V_P\rightarrow T_{DCS},
\end{equation}
其中$T_{DCS}$为实际DCS测点集合。一个领域参数可对应一个或多个冗余测点，也可根据具体诊断任务映射到筛选后的状态参数集合。以VHP工程算例为例，第二章状态监测模型首先预测27个候选输出，结果评价及后续故障诊断采用剔除$10MAA50CT132A$、$10MAA50CT151A$和$10MAA50CT152A$后的24个状态参数。下一章将利用$\mathcal M$从全系统知识图谱中抽取与这些状态参数相关的参数子图，并把状态监测模型得到的动态状态偏差作为图节点特征。

当前107节点/114三元组用于形成可复现的核心语义层和方法验证基础。随着运行日志、缺陷记录和检修文本持续接入，同一schema可继续扩展实体和关系，而无需改变后续KG-GCN的基本输入形式。

\subsection{本章小结}

本章针对汽轮机系统故障知识分散于系统结构、DCS测点、运行检修资料和故障文献中的特点，构建了面向故障诊断的领域知识图谱。首先依据汽轮机主通流、再热、抽汽回热和冷端关系分析了典型热力故障的传播特征，明确故障异常通常表现为具有设备连接关系的多参数协同变化；随后定义设备、参数、异常、故障和处置五类实体及相应关系，形成候选三元组抽取、规则校验、实体标准化和证据审核的知识构建流程；最终获得包含107个节点和114条三元组的核心语义层，并建立领域参数与实际DCS测点之间的映射。上述知识图谱为下一章从系统知识中提取参数关联拓扑、融合动态状态偏差并构建KG-GCN故障诊断模型提供了静态知识基础。
